\begin{abstract}
Removing 5\% of vision tokens collapses SmolVLA from 59.2\% to 0.0\% success rate, while the same reduction on $\pi_0.5$ retains 98.2\%. Token pruning, well-established for Vision Transformers, transfers unpredictably across Vision-Language-Action (VLA) models because the vision connector governs token compressibility. We analyze this sensitivity along two axes: connector coupling (token-independent vs.\ token-entangling) and task complexity (single-step vs.\ multi-step manipulation). The coefficient of variation (CV) of post-connector token L2 norms, ranging from 0.038 to 0.314 across models, serves as a single-pass diagnostic that correlates with both pruning tolerance and method sensitivity. Token-entangling connectors exhibit task-amplified collapse: X-VLA retains its full 99.2\% baseline on single-step Object manipulation at 90\% retention but loses 27.8 percentage points on multi-step Goal tasks at the same ratio. Low-CV models show 15$\times$ greater sensitivity to pruning method choice than higher-CV models, and at aggressive compression random selection outperforms L2 norm ranking by 36 percentage points for low-CV architectures. For pixel-shuffle connectors where all pruning strategies yield 0.0\%, Vision Token Re-encoder (VTR), a lightweight cross-attention module, recovers over 60\% success rate at 75\% compression through learned re-encoding. Evaluation spans six VLAs (450M to 7B parameters), three connector families, and four LIBERO task suites.
\end{abstract}
